\documentclass{article}
\usepackage[spanish]{babel}
\usepackage{multicol}
\usepackage{amsmath}
\usepackage{amsfonts}
\usepackage{enumerate}
\usepackage[utf8]{inputenc}
\usepackage{amssymb}
\usepackage{graphicx}
\usepackage{wasysym}
\usepackage{moodle}
\usepackage[usenames]{color}
\newcommand{\radio}{\ooalign{\hidewidth$\bullet$\hidewidth\cr$\ocircle$}}
\newcommand{\R}{$\mathbb{R}^3$ }
\newcommand{\wi}{\displaystyle \int}
\newtheorem{theorem}{Theorem}
\usepackage[left=2.00cm, right=2.00cm, top=2.00cm, bottom=2.00cm]{geometry}
\title{Parcial 2}
\author{Departamento de matemáticas}
\date{\today}

\begin{document}
	\maketitle
	\begin{quiz}{}
    \begin{cloze}{}
    	Considere la EDP parab\'olica
    	\[\begin{cases}
    		\displaystyle\frac{\partial u}{\partial t} = 4\frac{\partial^2 u}{\partial x^2}, \ x\in(0,1), \ t\in[0,1] \\
    		u(0,t) = e^{-t}, \ u(1,t)=\sin(t) \\
    		u(x,0) = x(1-x)
    	\end{cases}\]
       \begin{numerical}[points=15]{}
       	Se desea emplear el m\'etodo de diferencias finitas progresivo  (adelante en tiempo) para aproximar la soluci\'on del problema, si $h=\frac{1}{10}$ entonces el m\'aximo valor $k$ (con cinco cifras) que se debe tomar para que el esquema sea estable es $k=$\item 0.00125
       	\end{numerical}
        \begin{numerical}[points=18]{}
        	Corriendo el m\'etodo con los valores de $h=\frac{1}{10}$ y $k$ anterior se obtiene un valor aproximado (con tres cifras) para $u(0.1,\ 0.125)$ de \item[tolerance=0.003] 0.808
        \end{numerical}
    \end{cloze}
    \begin{cloze}{}
	Considere la EDP parab\'olica
	\[\begin{cases}
		\displaystyle\frac{\partial u}{\partial t} = 4\frac{\partial^2 u}{\partial x^2}, \ x\in(0,1), \ t\in[0,1] \\
		u(0,t) = e^{-t}, \ u(1,t)=\sin(t) \\
		u(x,0) = x(1-x)
	\end{cases}\]
	\begin{numerical}[points=15]{}
		Se desea emplear el m\'etodo de diferencias finitas progresivo  (adelante en tiempo) para aproximar la soluci\'on del problema, si $h=\frac{1}{10}$ entonces el m\'aximo valor $k$ (con cinco cifras) que se debe tomar para que el esquema sea estable es $k=$\item 0.00125
	\end{numerical}
	\begin{numerical}[points=18]{}
		Corriendo el m\'etodo con los valores de $h=\frac{1}{10}$ y $k$ anterior se obtiene un valor aproximado (con tres cifras) para $u(0.2,\ 0.125)$ de \item[tolerance=0.003] 0.732
	\end{numerical}
\end{cloze}
    \begin{cloze}{}
	Considere la EDP parab\'olica
	\[\begin{cases}
		\displaystyle\frac{\partial u}{\partial t} = 4\frac{\partial^2 u}{\partial x^2}, \ x\in(0,1), \ t\in[0,1] \\
		u(0,t) = e^{-t}, \ u(1,t)=\sin(t) \\
		u(x,0) = x(1-x)
	\end{cases}\]
	\begin{numerical}[points=15]{}
		Se desea emplear el m\'etodo de diferencias finitas progresivo  (adelante en tiempo) para aproximar la soluci\'on del problema, si $h=\frac{1}{10}$ entonces el m\'aximo valor $k$ (con cinco cifras) que se debe tomar para que el esquema sea estable es $k=$\item 0.00125
	\end{numerical}
	\begin{numerical}[points=18]{}
		Corriendo el m\'etodo con los valores de $h=\frac{1}{10}$ y $k$ anterior se obtiene un valor aproximado (con tres cifras) para $u(0.3,\ 0.125)$ de \item[tolerance=0.003] 0.655
	\end{numerical}
\end{cloze}
    \begin{cloze}{}
	Considere la EDP parab\'olica
	\[\begin{cases}
		\displaystyle\frac{\partial u}{\partial t} = 4\frac{\partial^2 u}{\partial x^2}, \ x\in(0,1), \ t\in[0,1] \\
		u(0,t) = e^{-t}, \ u(1,t)=\sin(t) \\
		u(x,0) = x(1-x)
	\end{cases}\]
	\begin{numerical}[points=15]{}
		Se desea emplear el m\'etodo de diferencias finitas progresivo  (adelante en tiempo) para aproximar la soluci\'on del problema, si $h=\frac{1}{10}$ entonces el m\'aximo valor $k$ (con cinco cifras) que se debe tomar para que el esquema sea estable es $k=$\item 0.00125
	\end{numerical}
	\begin{numerical}[points=18]{}
		Corriendo el m\'etodo con los valores de $h=\frac{1}{10}$ y $k$ anterior se obtiene un valor aproximado (con tres cifras) para $u(0.4,\ 0.125)$ de \item[tolerance=0.003] 0.578
	\end{numerical}
\end{cloze}
    \begin{cloze}{}
	Considere la EDP parab\'olica
	\[\begin{cases}
		\displaystyle\frac{\partial u}{\partial t} = 4\frac{\partial^2 u}{\partial x^2}, \ x\in(0,1), \ t\in[0,1] \\
		u(0,t) = e^{-t}, \ u(1,t)=\sin(t) \\
		u(x,0) = x(1-x)
	\end{cases}\]
	\begin{numerical}[points=15]{}
		Se desea emplear el m\'etodo de diferencias finitas progresivo  (adelante en tiempo) para aproximar la soluci\'on del problema, si $h=\frac{1}{10}$ entonces el m\'aximo valor $k$ (con cinco cifras) que se debe tomar para que el esquema sea estable es $k=$\item 0.00125
	\end{numerical}
	\begin{numerical}[points=18]{}
		Corriendo el m\'etodo con los valores de $h=\frac{1}{10}$ y $k$ anterior se obtiene un valor aproximado (con tres cifras) para $u(0.5,\ 0.125)$ de \item[tolerance=0.003] 0.5
	\end{numerical}
\end{cloze}
    \begin{cloze}{}
	Considere la EDP parab\'olica
	\[\begin{cases}
		\displaystyle\frac{\partial u}{\partial t} = 4\frac{\partial^2 u}{\partial x^2}, \ x\in(0,1), \ t\in[0,1] \\
		u(0,t) = e^{-t}, \ u(1,t)=\sin(t) \\
		u(x,0) = x(1-x)
	\end{cases}\]
	\begin{numerical}[points=15]{}
		Se desea emplear el m\'etodo de diferencias finitas progresivo  (adelante en tiempo) para aproximar la soluci\'on del problema, si $h=\frac{1}{10}$ entonces el m\'aximo valor $k$ (con cinco cifras) que se debe tomar para que el esquema sea estable es $k=$\item 0.00125
	\end{numerical}
	\begin{numerical}[points=18]{}
		Corriendo el m\'etodo con los valores de $h=\frac{1}{10}$ y $k$ anterior se obtiene un valor aproximado (con tres cifras) para $u(0.6,\ 0.125)$ de \item[tolerance=0.003] 0.422
	\end{numerical}
\end{cloze}
    \begin{cloze}{}
	Considere la EDP parab\'olica
	\[\begin{cases}
		\displaystyle\frac{\partial u}{\partial t} = 4\frac{\partial^2 u}{\partial x^2}, \ x\in(0,1), \ t\in[0,1] \\
		u(0,t) = e^{-t}, \ u(1,t)=\sin(t) \\
		u(x,0) = x(1-x)
	\end{cases}\]
	\begin{numerical}[points=15]{}
		Se desea emplear el m\'etodo de diferencias finitas progresivo  (adelante en tiempo) para aproximar la soluci\'on del problema, si $h=\frac{1}{10}$ entonces el m\'aximo valor $k$ (con cinco cifras) que se debe tomar para que el esquema sea estable es $k=$\item 0.00125
	\end{numerical}
	\begin{numerical}[points=18]{}
		Corriendo el m\'etodo con los valores de $h=\frac{1}{10}$ y $k$ anterior se obtiene un valor aproximado (con tres cifras) para $u(0.7,\ 0.125)$ de \item[tolerance=0.003] 0.345
	\end{numerical}
\end{cloze}
    \end{quiz}
\end{document}